\newcommand{\drawHaplotypePresentation}[1]{%
  \begin{tikzpicture}[scale=0.74, x=1.0cm, y=1.2cm, font=\footnotesize]
    % 1. CHROMOSOME LINES (GROUND TRUTH)
    \draw[faint, dashed] (-1.5, 3.5) -- (13.5, 3.5) node[right, text=black, align=left] {\textbf{Haplotype 1}\\(vérité {\Large $\boldsymbol{\Sigma}$} $= +$)};
    \draw[faint, dashed] (-1.5, 0) -- (13.5, 0) node[right, text=black, align=left] {\textbf{Haplotype 2}\\(vérité {\Large $\boldsymbol{\Sigma}$} $= -$)};

    % 2. BACKGROUND: READS AND SEGMENTS (step >= 2)
    \ifnum#1>1\relax
      \begin{scope}[on background layer]
        % --- Reads (Segments) ---
        \draw[readlink] (1, 3.5) -- (1, 4.0); \draw[read] (-1.2, 4.0) -- (3.2, 4.0); 
        \draw[readlink] (3, 3.5) -- (3, 4.4); \draw[read] (0.8, 4.4) -- (5.2, 4.4); 
        \draw[readlink] (5, 3.5) -- (5, 4.8); \draw[read] (2.8, 4.8) -- (7.2, 4.8); 
        \draw[readlink] (7, 3.5) -- (7, 4.0); \draw[read] (4.8, 4.0) -- (9.2, 4.0); 
        \draw[readlink] (9, 3.5) -- (9, 4.4); \draw[read] (6.8, 4.4) -- (11.2, 4.4); 
        \draw[readlink] (11, 3.5) -- (11, 4.8); \draw[read] (8.8, 4.8) -- (13.2, 4.8); 

        % Haplotype 2
        \draw[readlink] (2, 0) -- (2, 0.5); \draw[read] (-0.2, 0.5) -- (4.2, 0.5); 
        \draw[readlink] (4, 0) -- (4, 0.9); \draw[read] (1.8, 0.9) -- (6.2, 0.9); 
        \draw[readlink] (6, 0) -- (6, 1.3); \draw[read] (3.8, 1.3) -- (8.2, 1.3); 
        \draw[readlink] (8, 0) -- (8, 0.5); \draw[read] (5.8, 0.5) -- (10.2, 0.5); 
        \draw[readlink] (10, 0) -- (10, 0.9); \draw[read] (7.8, 0.9) -- (12.2, 0.9); 
        \draw[readlink] (12, 0) -- (12, 1.3); \draw[read] (9.8, 1.3) -- (14.2, 1.3); 
      \end{scope}

      % NODES (SPIN CENTERS) - Now visible starting from step 2!
      \node[spin, fill=gray!8] (t1) at (1, 3.5) {$+$};
      \node[spin, fill=gray!8] (t2) at (3, 3.5) {$+$};
      \node[spin, fill=gray!8] (t3) at (5, 3.5) {$+$};
      \node[spin, fill=gray!8] (t4) at (7, 3.5) {$+$};
      \node[spin, fill=gray!8] (t5) at (9, 3.5) {$+$};
      \node[spin, fill=gray!8] (t6) at (11, 3.5) {$+$};

      \node[spin, fill=gray!20] (b1) at (2, 0) {$-$};
      \node[spin, fill=gray!20] (b2) at (4, 0) {$-$};
      \node[spin, fill=gray!20] (b3) at (6, 0) {$-$};
      \node[spin, fill=gray!20] (b4) at (8, 0) {$-$};
      \node[spin, fill=gray!20] (b5) at (10, 0) {$-$};
      \node[spin, fill=gray!20] (b6) at (12, 0) {$-$};
    \fi

    % 3. SIGNED GRAPH EDGES ONLY (step >= 3)
    \ifnum#1>2\relax
      \begin{scope}[on background layer]
        % --- Inter-Chromosome Edges (True Negatives) ---
        \draw[interedge] (1, 3.5) -- (2, 0); 
        \draw[interedge] (3, 3.5) -- (4, 0); 
        \draw[interedge] (3, 3.5) -- (6, 0); 
        \draw[interedge] (5, 3.5) -- (2, 0); 
        \draw[interedge] (5, 3.5) -- (4, 0); 
        \draw[interedge] (5, 3.5) -- (6, 0); 
        \draw[interedge] (7, 3.5) -- (4, 0); 
        \draw[interedge] (7, 3.5) -- (6, 0); 
        \draw[interedge] (7, 3.5) -- (8, 0); 
        \draw[interedge] (7, 3.5) -- (10, 0); 
        \draw[interedge] (9, 3.5) -- (8, 0); 
        \draw[interedge] (9, 3.5) -- (10, 0); 
        \draw[interedge] (9, 3.5) -- (12, 0); 
        \draw[interedge] (11, 3.5) -- (8, 0); 
        \draw[interedge] (11, 3.5) -- (12, 0); 

        % False Positives (FP) => 25% noise
        \draw[posedge, draw=red!80!black] (1, 3.5) -- node[errlabel, pos=0.2] {FP} (4, 0); 
        \draw[posedge, draw=red!80!black] (3, 3.5) -- node[errlabel, pos=0.8] {FP} (2, 0); 
        \draw[posedge, draw=red!80!black] (5, 3.5) -- node[errlabel, pos=0.2] {FP} (8, 0); 
        \draw[posedge, draw=red!80!black] (9, 3.5) -- node[errlabel, pos=0.8] {FP} (6, 0); 
        \draw[posedge, draw=red!80!black] (11, 3.5) -- node[errlabel, pos=0.2] {FP} (10, 0); 
      \end{scope}

      % INTRA-EDGES
      \draw[posedge] (t1) to[bend right=25] (t2);
      \draw[negedge, draw=red!80!black] (t2) to[bend right=25] node[errlabel] {FN} (t3);
      \draw[posedge] (t3) to[bend right=25] (t4);
      \draw[posedge] (t4) to[bend right=25] (t5);
      \draw[posedge] (t5) to[bend right=25] (t6);
      
      \draw[posedge] (t1) to[bend right=35] (t3);
      \draw[posedge] (t2) to[bend right=35] (t4);
      \draw[negedge, draw=red!80!black] (t3) to[bend right=35] node[errlabel] {FN} (t5);
      \draw[posedge] (t4) to[bend right=35] (t6);

      \draw[posedge] (b1) to[bend right=25] (b2);
      \draw[posedge] (b2) to[bend right=25] (b3);
      \draw[posedge] (b3) to[bend right=25] (b4);
      \draw[negedge, draw=red!80!black] (b4) to[bend right=25] node[errlabel] {FN} (b5);
      \draw[posedge] (b5) to[bend right=25] (b6);
      
      \draw[negedge, draw=red!80!black] (b1) to[bend right=35] node[errlabel] {FN} (b3);
      \draw[negedge, draw=red!80!black] (b2) to[bend right=35] node[errlabel] {FN} (b4);
      \draw[posedge] (b3) to[bend right=35] (b5);
      \draw[posedge] (b4) to[bend right=35] (b6);
    \fi

    % 5. LEGEND
    \ifnum#1=2\relax
      \node[scale=0.65,chapbox, align=left, fill=white] at (0,1.75){
        \begin{tabular}{ll}
          \tikz[baseline=-0.5ex]{\draw[line width=5pt, gray!45, line cap=round] (0,0) -- (0.6,0);} & Lecture ADN (segment) \\
          \tikz[baseline=-0.5ex]{\node[spin, scale=0.6, fill=gray!8] at (0,0) {$+$};} & Nœud du graphe (vérité {\Large $\boldsymbol{\Sigma}$})
        \end{tabular}
      };
    \fi
    \ifnum#1>2\relax
      \node[scale=0.65,chapbox, align=left, fill=white] at (0,1.75){
        \begin{tabular}{ll}
          \tikz[baseline=-0.5ex]{\draw[line width=5pt, gray!45, line cap=round] (0,0) -- (0.6,0);} & Lecture ADN (nœud) \\
          \tikz[baseline=-0.5ex]{\draw[posedge] (0,0) -- (0.6,0);} & Attraction (trait plein)  \\
          \tikz[baseline=-0.5ex]{\draw[negedge] (0,0) -- (0.6,0);} & Répulsion (trait pointillé) \\
          \textcolor{red!80!black}{\textbf{FP}} & Faux positif (bruit)  \\
          \textcolor{red!80!black}{\textbf{FN}} & Faux négatif (bruit)
        \end{tabular}
      };
    \fi
  \end{tikzpicture}
}
